% !TEX root = ../main.tex
\chapter{Contesto, obiettivi e strumenti}
\label{cap:contesto}

\section{Contesto applicativo e obiettivi}
Venue Event Manager è un progetto universitario dedicato alla gestione di strutture per eventi,
delle loro risorse e delle attività di organizzazione e prenotazione.
Il sistema distingue la presentazione di richieste di evento dalla gestione degli eventi programmati
e comprende utenti, liste di invitati, prenotazioni, biglietti, recensioni e report amministrativi.
L'obiettivo del lavoro è sviluppare e documentare un nucleo applicativo che renda esplicite
le regole del dominio e ne gestisca la persistenza in una base di dati relazionale.

\section{Perimetro della versione realizzata}
La versione descritta comprende modelli di dominio, servizi applicativi e repository JDBC
per PostgreSQL.
Le funzionalità sono esercitabili attraverso i servizi e i test automatici.
Non sono attualmente presenti un'interfaccia grafica, un'API web o una CLI definitiva.
La verifica delle credenziali non crea una sessione applicativa e lo stato di attesa di pagamento
di una prenotazione non corrisponde all'integrazione con un sistema di pagamento reale.
Questi confini guidano l'analisi dei requisiti e la discussione dei risultati.

\section{Tecnologie e strumenti}
\bozza{Descrivere ruolo e motivazione di Java, Maven, PostgreSQL, JDBC, Flyway,
HikariCP, JUnit, Mockito e JaCoCo; verificare versioni nel pom.xml.
Documentare editor, strumenti per diagrammi e ambiente effettivamente usati dal gruppo.}

\section{Organizzazione della relazione}
\bozza{Completare una volta stabilizzato l'indice, richiamando lo scopo dei capitoli successivi.}
